\section{Limitations, Roadmap and Conclusion}\label{sec:conclusion}

\subsection{Limitations}
\begin{itemize}
  \item \textbf{Design-only components.} The EVPN-VXLAN fabric, SD-WAN, ZTNA/SASE, the wireless controller
        and 802.1X are specified but not implemented, because the lab images do not support them.
  \item \textbf{Lab substitutions} (\cref{tab:substitutions}): the zone-based firewall stands in for an NGFW,
        an IOS router for the DHCP/DNS/NTP servers, and 1\,Gbit/s virtual links for the 10/40\,Gbit/s
        backbone, whose throughput could therefore not be measured.
  \item \textbf{Remaining single point of failure.} Each branch has one access switch in the lab; its failure
        isolates that branch's wired hosts. Production branches need a second switch or stacked switches.
  \item \textbf{Stateless segmentation.} The VLAN ACLs admit return traffic by pattern, not by session state.
  \item \textbf{Simplified security credentials.} IKEv1 with pre-shared keys and local AAA accounts are lab
        simplifications.
  \item \textbf{Double failures.} Losing the HQ edge-to-edge link \emph{and} a branch's primary tunnel at the
        same time re-creates the routing loop described in \cref{sec:challenges}.
  \item \todo{commands the lab images rejected, and anything not verified in the lab}
\end{itemize}

\subsection{Roadmap}
\begin{enumerate}
  \item \textbf{BFD} on the tunnels and the provider links, to detect WAN failures in under a second instead of 9\,s.
  \item \textbf{IKEv2 with certificates}, \textbf{TACACS+/RADIUS} for device access and \textbf{802.1X} at the
        access layer, which also enables dynamic VLAN assignment.
  \item A \textbf{stateful NGFW pair} between internal zones and at the Internet edge.
  \item A second edge-to-edge link in a port-channel, removing the double-failure case.
  \item The \textbf{EVPN-VXLAN data-centre fabric} in the reserved block \texttt{10.1.128.0/18}.
  \item \textbf{SD-WAN} (or DMVPN as an intermediate step) and \textbf{SASE} for remote users and branch
        Internet traffic.
\end{enumerate}

\subsection{Conclusion}
The design answers TVI's four needs. High availability comes from paired devices, providers and tunnels
whose failover is driven by tracked routes and fast protocol timers. Branch traffic to the cloud leaves
locally, while inter-site traffic is encrypted and routed with three summarised prefixes. Guests, IoT devices
and corporate users are separated at the first hop and at the edge. The whole network is generated from one
data model, verified automatically before deployment and operated with scripts that back up every device
before they change it.
\todo{close with the measured results from \cref{tab:failover}}
